\chapter*{Conclusion}
\phantomsection
\addcontentsline{toc}{chapter}{Conclusion}

This internship addressed two complementary aspects of open-source power
electronics development: the validation of existing platforms and the
preparation of future high-voltage converter modules.

The first part focused on the SPIN Oscilloscope Shield. The SPOC board was
designed to route selected SPIN signals to oscilloscope channels and to support
automated hardware validation. The tests confirmed that the main routing paths
operate correctly through the multiplexer network and that the shield can be
used as a practical measurement interface.

The software developed around the SPOC board completes this validation platform. The \texttt{Core-CUT} firmware generates configurable test signals, the \texttt{Core-MUX} firmware selects the measurement paths, and the Python host software coordinates the oscilloscope and the validation procedure. This architecture makes the tests faster, more repeatable, and less dependent on manual oscilloscope manipulation. It also demonstrated its usefulness by identifying and analysing a burst-mode timing issue in the tested PWM generation framework.

The second part focused on the design of a GaN-based Half-Bridge Board. The
schematic and PCB routing were completed around a modular power-stage
architecture including digital isolation, isolated auxiliary supplies, gate
drivers, GaN transistors, local DC-link capacitors, and clearly separated
reference domains. This board prepares future work on higher-voltage and
higher-power applications, where layout, insulation, gate driving, thermal
behaviour, and electromagnetic disturbances become critical.

The substrate analysis showed that aluminium IMS improves thermal performance,
but also introduces manufacturing and end-of-life constraints.

Overall, the internship contributed to OwnTech's objective of building power
electronics hardware that is more open, modular, testable, and maintainable.
The SPOC platform strengthens the validation workflow of existing boards, while
the GaN Half-Bridge Board provides a basis for future experimental work. The
next steps are the design of the half-bridge control board and progressive
laboratory validation of the GaN power stage, first at safe low voltage and then
under high-voltage operating conditions.